\documentclass[11pt]{article}
% Tables of Hilbert functions for "Hilbert functions of filtered graphical
% zonotopal algebras" (B. Shapiro, I. Smirnov, A. Vaintrob).
% Formerly Appendix B of the paper.  Compile with pdflatex appendix_B.tex.
% The tables themselves are in tables345_clean.tex, the file that
% check_tables.py reads; do not edit one without the other.
\usepackage[margin=.8in]{geometry}
\usepackage{mathtools,amssymb}
\usepackage{graphicx,xcolor}
\usepackage{rotating}
\usepackage{booktabs,array}
\usepackage{tikz}
\tikzset{vertex/.style={circle,fill,inner sep=1.pt}}
\usepackage[mathscr]{eucal}
\usepackage{caption}
\usepackage{url}

\newcommand{\CC}{\mathscr{C}}
\newcommand{\hs}{\mathscr{H}}

% Statement numbers of the paper (first arXiv version) cited in the
% table source.
\makeatletter
\newcommand{\paperref}[2]{\expandafter\gdef\csname r@#1\endcsname{{#2}{}}}
\makeatother
\paperref{rem:small-complete-profiles}{6.13}
\paperref{thm:D4-strata}{6.10}

\renewcommand{\thesection}{B}
\renewcommand{\thetable}{B\arabic{table}}

\title{Tables of Hilbert functions of filtered graphical zonotopal algebras}
\author{Boris Shapiro, Ilya Smirnov, Arkady Vaintrob}
\date{Supplement to the paper \emph{Hilbert functions of filtered graphical
zonotopal algebras}}

\begin{document}
\maketitle

\section{Selected numerical tables}

These tables give the Hilbert functions $\hs_G^f$ of the filtered external
zonotopal algebras $\CC^f_G$ for a number of graphs $G$ and nondegenerate
polynomials $f$. Statement and equation numbers refer to the first arXiv
version of the paper; Tables B1--B4 were Tables 12--15 of its version of
21 October 2026. All entries have been
verified by exact rational arithmetic directly in the square-free edge
algebra $\Phi_G$, using the generators $Y_v=f(X_v)$: the script
\url{check_tables.py} in this directory reads the edge lists from the
source of the drawings below and recomputes every entry. The file
\url{tables.json} (and \url{tables.csv}) lists every entry with its edge
list.

If $f=u$, i.e.\ when $\CC^f_G$ is the usual graded zonotopal algebra, the
answer is given by Theorem~2.10. In particular, if $G$ is a tree with $n$
edges, then $\hs_G^u=(1+t)^n$. We assume that $\deg f\le \Delta(G)$, since,
as explained in Remark~3.3, the terms of $f$ of degree higher than
$\Delta(G)$ do not affect $\CC^f_G$.

For a nonzero quadratic coefficient and a graph whose loopless part has no
isolated vertices or isolated edges, Proposition~3.13(iii) gives
$h_1=|V|$. For graphs of maximum degree three the tables give complete
profiles; for maximum degree four they give values at nine parameters.

\subsection{Graphs with at most five vertices}

The tables use the actual generating function $Y_v=f(X_v)$. In
particular, the parameter $f$ is not the inverse series used in the
defining relations; compare (6.5). The graphs are determined by their
drawings: each drawn segment or arc is one edge.

% Clean version, 19 October 2026: the Delta=3 tables give complete cubic profiles
% (from clean/checks/complete_profiles.py); captions and guidance revised.
\newcommand{\Gafa}[1]{
  % D4 with first doubled edge =--
  \begin{tikzpicture}[scale=#1*.8]
    \node[vertex] (a) at (-0.1,0) {};
    \node[vertex] (b) at (0.8,0) {};
    \node[vertex] (c) at (1.6,0) {};
    \node[vertex] (d) at (2.4,0) {};
    \node (e) at (0,-.5) {};   % empty edge
    \draw (a) to[out=25,in=155, min distance=.7em] (b);
    \draw (a) to[out=-25,in=-155, min distance=.7em] (b);
    \draw (b)--(c)--(d);
  \end{tikzpicture}}

\newcommand{\Gafb}[1]{
  % D4 with middle doubled edge -=-
  \begin{tikzpicture}[scale=#1*.8]
    \node[vertex] (a) at (-0.1,0) {};
    \node[vertex] (b) at (0.7,0) {};
    \node[vertex] (c) at (1.6,0) {};
    \node[vertex] (d) at (2.4,0) {};
    \node (e) at (0,-.5) {};   % empty edge
    \draw (b) to[out=25,in=155, min distance=.7em] (c);
    \draw (b) to[out=-25,in=-155, min distance=.7em] (c);
    \draw (b)--(a) (c)--(d);
  \end{tikzpicture}}

\newcommand{\Gafc}[1]{
  % A4 with first and third doubled edges =-=
  \begin{tikzpicture}[scale=#1*.8]
    \node[vertex] (a) at (-0.1,0) {};
    \node[vertex] (b) at (0.8,0) {};
    \node[vertex] (c) at (1.6,0) {};
    \node[vertex] (d) at (2.5,0) {};
    \node (e) at (0,-.5) {};   % empty edge
    \draw (a) to[out=25,in=155, min distance=.7em] (b);
    \draw (a) to[out=-25,in=-155, min distance=.7em] (b);
    \draw (c) to[out=25,in=155, min distance=.7em] (d);
    \draw (c) to[out=-25,in=-155, min distance=.7em] (d);
    \draw (b)--(c);
  \end{tikzpicture}}


\newcommand{\Gppw}[1]{
  % Paw graph w doubled edge -<||
  \begin{tikzpicture}[scale=#1]
    \node[vertex] (d) at (0,0) {};
    \node[vertex] (a) at (30:.8) {};
    \node[vertex] (b) at (180:.8) {};
    \node[vertex] (c) at (-30:.8) {};
    \draw (b)--(d)--(c) (a)--(d);
    \draw (a) to[out=-65,in=65, min distance=.6em] (c);
    \draw (a) to[out=-105,in=115, min distance=.6em] (c);
\end{tikzpicture}}


\newcommand{\Gcfa}[1]{
  % 4-cycle with a doubled edge  <\geq
  \begin{tikzpicture}[scale=#1*.8]
    \node[vertex] (a) at (0:.8) {};
    \node[vertex] (b) at (90:.8) {};
    \node[vertex] (c) at (180:.8) {};
    \node[vertex] (d) at (-90:.8) {};
    \node (f) at (90:1.) {};    % empty vertex
    \draw (b)--(c)--(d)--(a);
    \draw (a) to[out=120,in=-30, min distance=1.1em] (b);
    \draw (a) to[out=150,in=-60, min distance=1.1em] (b);
  \end{tikzpicture}}

\newcommand{\Gcfb}[1]{
  % 4-cycle with two doubled edges  <\geq
  \begin{tikzpicture}[scale=#1*.8]
    \node[vertex] (a) at (0:.8) {};
    \node[vertex] (b) at (90:.8) {};
    \node[vertex] (c) at (180:.8) {};
    \node[vertex] (d) at (-90:.8) {};
%   \node (f) at (90:1) {};    % empty vertex
%    \node (g) at (180:1) {};    % empty vertex
    \draw (b)--(c) (d)--(a);
    \draw (a) to[out=120,in=-30, min distance=1.1em] (b);
    \draw (a) to[out=150,in=-60, min distance=1.1em] (b);
    \draw (d) to[out=120,in=-30, min distance=1.1em] (c);
    \draw (d) to[out=150,in=-60, min distance=1.1em] (c);
  \end{tikzpicture}}

\newcommand{\GDfo}[1]{
  % D4 tree
  \begin{tikzpicture}[scale=#1*.8]
    \node[vertex] (d) at (0,0) {};
    \node[vertex] (a) at (60:.8) {};
    \node[vertex] (b) at (180:.8) {};
    \node[vertex] (c) at (-60:.8) {};
    \draw (a)--(d)--(c);
    \draw (d)--(b);
  \end{tikzpicture}}

\newcommand{\GPaw}[1]{
  % Paw graph  -<|
  \begin{tikzpicture}[scale=#1]
    \node[vertex] (d) at (0,0) {};
    \node[vertex] (a) at (30:.8) {};
    \node[vertex] (b) at (180:.8) {};
    \node[vertex] (c) at (-30:.8) {};
    \draw (b)--(d)--(c)--(a)--(d);
  \end{tikzpicture}}

\newcommand{\GDia}[1]{
  % Diamond graph = K4-e <|>
  \begin{tikzpicture}[scale=#1*.8]
    \node[vertex] (a) at (0:.8) {};
    \node[vertex] (b) at (90:.8) {};
    \node[vertex] (c) at (180:.8) {};
    \node[vertex] (d) at (-90:.8) {};
    \node (f) at (90:1) {};    % empty vertex
    \draw (a)--(b)--(c)--(d)--(a)--(c);
  \end{tikzpicture}}

\newcommand{\GKfo}[1]{
  % Complete graph K4
  \begin{tikzpicture}[scale=#1*.8]
    \node[vertex] (d) at (0,0) {};
    \node[vertex] (a) at (60:.8) {};
    \node[vertex] (b) at (180:.8) {};
    \node[vertex] (c) at (-60:.8) {};
    \draw (a)--(b)--(c)--(a)--(d)--(c);
    \draw (d)--(b);
    % \node[vertex] (o) at (0:0) {};
    % \node[vertex] (b) at (90:.8) {};
    % \node[vertex] (c) at (180:.8) {};
    % \node[vertex] (d) at (-90:.8) {};
    % \draw (a)--(b)--(c)--(d)--(a)--(c) (b)--(d);
  \end{tikzpicture}}

\newcommand{\GDfi}[1]{
  % D5 tree
  \begin{tikzpicture}[scale=#1*.8]
  \node[vertex] (c) at (0,0) {};
  \node[vertex] (a) at (-.7,0) {};
  \node[vertex] (b) at (-1.4,0) {};
  \node[vertex] (r1) at (1.0,.6) {};
  \node[vertex] (r2) at (1.0,-.6) {};
  \draw (a)--(b)--(c);
  \draw (r2)--(c)--(r1);
\end{tikzpicture}}

\newcommand{\Glol}[1]{
% Lollipop L_{3,2}  (K3 + path of length 2 attached at one vertex)
  \begin{tikzpicture}[scale=#1*.8]
  \node[vertex] (c) at (0,0) {};
  \node[vertex] (a) at (-.7,0) {};
  \node[vertex] (b) at (-1.4,0) {};
  \node[vertex] (r1) at (1.0,.6) {};
  \node[vertex] (r2) at (1.0,-.6) {};
  \draw (a)--(b)--(c);
  \draw (r1)--(r2)--(c)--(r1);
\end{tikzpicture}
}

\newcommand{\Gban}[1]{
  \begin{tikzpicture}[scale=#1*.8]
% Banner Graph (C4 with a pendant edge)
\node[vertex] (v1) at (0,0) {};
\node[vertex] (v2) at (1.4,0) {};
\node[vertex] (v3) at (.7,.7) {};
\node[vertex] (v4) at (.7,-.7) {};
\node[vertex] (p)  at (-.75,0) {};
\draw (v1)--(v4)--(v2)--(v3)--(v1);
\draw (v1)--(p);
\end{tikzpicture}
}

\newcommand{\Gbul}[1]{
%  Bull Graph (triangle with two pendant edges on distinct vertices)
  \begin{tikzpicture}[scale=#1*.8]
  \node[vertex] (c) at (0,0) {};
  \node[vertex] (h1) at (1.7,.6) {};
  \node[vertex] (h2) at (1.7,-.6) {};
  \node[vertex] (r1) at (1.0,.6) {};
  \node[vertex] (r2) at (1.0,-.6) {};
  \draw (h1)--(r1) (h2)--(r2);
  \draw (r1)--(r2)--(c)--(r1);
\end{tikzpicture}
}

\newcommand{\Ghou}[1]{
  % House Graph
  \begin{tikzpicture}[scale=#1*.7]
  \node[vertex] (a) at (45:1) {};
  \node[vertex] (b) at (135:1) {};
  \node[vertex] (c) at (-135:1) {};
  \node[vertex] (d) at (-45:1) {};
  \node[vertex] (e) at (90:1.5) {};
  \node (f) at (90:1.8) {};  % empty vertex
%
  \draw (a)--(b)--(c)--(d)--(a); % square
  \draw (a)--(e)--(b);          % roof
\end{tikzpicture}
}

\newcommand{\Gkit}[1]{
% Kite Graph (Diamond K4-e with a pendant vertex)
  \begin{tikzpicture}[scale=#1*.9]
  \node[vertex] (v1) at (0,0) {};
  \node[vertex] (v2) at (1.5,0) {};
  \node[vertex] (v3) at (.75,.75) {};
  \node[vertex] (v4) at (.75,-.75) {};
  \node[vertex] (p)  at (-.75,0) {};
  \draw (v2)--(v3)--(v1);
  \draw (v1)--(v4)--(v2)
  (v3)--(v4);
  \draw (v1)--(p);
  % \node at (.6,-1.45) {Kite graph};
\end{tikzpicture}
}

\newcommand{\Gbip}[1]{
% $K_{3,2}$
  \begin{tikzpicture}[scale=#1*.9]
  \node[vertex] (v1) at (0,0) {};
  \node[vertex] (v2) at (1.5,0) {};
  \node[vertex] (v3) at (.75,.75) {};
  \node[vertex] (v4) at (.75,-.75) {};
  \node[vertex] (c)  at (.75,0) {};
  \draw (v1)--(v3)--(v2)--(v4)--(v1);
  \draw (v1)--(c)--(v2);
\end{tikzpicture}
}

\newcommand{\Gsub}[1]{
% $K_4$ with a subdivided edge
  \begin{tikzpicture}[scale=#1*.9]
    \node[vertex] (v1) at (0,0) {};
    \node[vertex] (v2) at (1.5,0) {};
    \node[vertex] (v3) at (.75,.75) {};
    \node[vertex] (v4) at (.75,-.75) {};
    \node[vertex] (c)  at (.75,0) {};
    \draw (v1)--(v3)--(v2)--(v4)--(v1);
    \draw (v1)--(c)--(v2);
    \draw (v3)--(c);
\end{tikzpicture}
}

\newcommand{\Gsta}[1]{
% S5 star
  \begin{tikzpicture}[scale=#1]
  \node[vertex] (a) at (0,0) {};
  \node[vertex] (b) at (0.8,0) {};
  \node[vertex] (c) at (-0.8,0) {};
  \node[vertex] (d) at (0,0.8) {};
  \node[vertex] (e) at (0,-0.8) {};
  \draw (c)--(a)--(b);
  \draw (d)--(a)--(e);
%  \node at (0.1,-1.25) {S5 star};
\end{tikzpicture}
}

\newcommand{\Gcri}[1]{
% Cricket Graph
  \begin{tikzpicture}[scale=#1]
  \node[vertex] (c) at (0,0) {};
  \node[vertex] (l1) at (-.7,.4) {};
  \node[vertex] (l2) at (-.7,-.4) {};
  \node[vertex] (r1) at (1.0,.6) {};
  \node[vertex] (r2) at (1.0,-.6) {};
  \draw (l2)--(c)--(l1);
  \draw (r1)--(r2)--(c)--(r1);
\end{tikzpicture}
}

\newcommand{\Gbut}[1]{
  % Butterfly Graph |><|
  \begin{tikzpicture}[scale=#1]
    \node[vertex] (c) at (0,0) {};
    \node[vertex] (l1) at (-1.,.6) {};
    \node[vertex] (l2) at (-1.,-.6) {};
    \node[vertex] (r1) at (1.,.6) {};
    \node[vertex] (r2) at (1.,-.6) {};
    \draw (l1)--(l2)--(c)--(l1);
    \draw (r1)--(r2)--(c)--(r1);
    % \node at (0,-.95) {Butterfly Graph};
  \end{tikzpicture}
}

\newcommand{\Gdar}[1]{
  % Dart graph
  \begin{tikzpicture}[scale=#1]
    \node[vertex] (v1) at (0,0) {};
    \node[vertex] (v2) at (1.5,0) {};
    \node[vertex] (v3) at (.75,.75) {};
    \node[vertex] (v4) at (.75,-.75) {};
    \node[vertex] (p)  at (-.75,0) {};
    \draw (v1)--(v2)--(v3)--(v1);
    \draw (v1)--(v4)--(v2);
    \draw (v1)--(p);
    % \node at (.6,-1.45) {Dart graph};
  \end{tikzpicture}
}

\newcommand{\Ggem}[1]{
  % Gem graph (P_4 + K_1): a 4-vertex path plus a universal vertex
  \begin{tikzpicture}[scale=#1*1.1]
     \node[vertex] (v1) at (-36:1.2) {};
     \node[vertex] (v2) at (-72:1.2) {};
     \node[vertex] (v3) at (-108:1.2) {};
     \node[vertex] (v4) at (-144:1.2) {};
     \node[vertex] (u)  at (0,0) {};
    \draw (v1)--(v2)--(v3)--(v4);
    \draw (u)--(v1) (u)--(v2) (u)--(v3) (u)--(v4);
  \end{tikzpicture}
}

\newcommand{\GKKK}[1]{
%  K2+3K1
    \begin{tikzpicture}[scale=#1*.75]
      \node[vertex] (v1) at (90:1.1) {};
      \node[vertex] (v2) at (90:0.3) {};
      \node[vertex] (v3) at (-45:1.3) {};
      \node[vertex] (v4) at (-135:1.3) {};
      \node[vertex] (v5) at (90:-.4) {};
      \draw   (v1)--(v3) (v1)--(v4)
      (v2)--(v3) (v2)--(v4)
      (v3)--(v4) (v3)--(v5) (v4)--(v5);
      % \node at (0,-1.7) {$K_2+3K_1$};
    \end{tikzpicture}
  }

  \newcommand{\Gllo}[1]{
    % Lollipop L_{4,1}  (K4 with a single pendant vertex)
    \begin{tikzpicture}[scale=#1*.9]
    \node[vertex] (v1) at (0,0) {};
    \node[vertex] (v2) at (.75,0) {};
    \node[vertex] (v3) at (1.22,.75) {};
    \node[vertex] (v4) at (1.22,-.75) {};
    \node[vertex] (p)  at (-.75,0) {};
    \draw (v1)--(v2)--(v3)--(v1);
    \draw (v1)--(v4)--(v2) (v3)--(v4);
    \draw (v1)--(p);
  \end{tikzpicture}
}

\newcommand{\Gwhe}[1]{
  % Wheel graph W_5  (C4 plus a hub)
  \begin{tikzpicture}[scale=#1]
    \node[vertex] (v1) at (0,0) {};
    \node[vertex] (v2) at (1.5,0) {};
    \node[vertex] (v3) at (.75,.75) {};
    \node[vertex] (v4) at (.75,-.75) {};
    \node[vertex] (c)  at (.75,0) {};
    \draw (v1)--(v3)--(v2)--(v4)--(v1);
    \draw (v1)--(c)--(v3);
    \draw (v2)--(c)--(v4);
  \end{tikzpicture}
}

\newcommand{\Ghox}[1]{
  % House X Graph (house with both diagonals in the square)
  \begin{tikzpicture}[scale=#1*0.75]
    \node[vertex] (v1) at (90:1.1) {};
    \node[vertex] (v2) at (10:1.) {};
    \node[vertex] (v3) at (-90:1.1) {};
    \node[vertex] (v4) at (-90:.38) {};
    \node[vertex] (v5) at (170:1.) {};
    \draw (v1)--(v2)--(v3)--(v4)--(v5)--(v1)
    (v2)--(v4) (v2)--(v5)
    (v3)--(v5);
  \end{tikzpicture}
}

\newcommand{\GKme}[1]{
%  K5-e
  \begin{tikzpicture}[scale=#1*.75]
    \node[vertex] (v1) at (90:1.1) {};
    \node[vertex] (v2) at (90:0.3) {};
    \node[vertex] (v3) at (-45:1.3) {};
    \node[vertex] (v4) at (-135:1.3) {};
    \node[vertex] (v5) at (90:-.4) {};
    \draw (v1)--(v2)
    (v1)--(v3) (v1)--(v4) (v1)--(v5)
    (v2)--(v3) (v2)--(v4)
    (v3)--(v4) (v3)--(v5)
    (v4)--(v5);
    % \node at (0,-1.7) {$K_3+2K_1=K_{5}-e$};
  \end{tikzpicture}
}

\newcommand{\GKfi}[1]{
  % Complete graph K_5
  \begin{tikzpicture}[scale=#1*0.8]
    \node[vertex] (v1) at (90:1.) {};
    \node[vertex] (v2) at (18:1.) {};
    \node[vertex] (v3) at (-54:1.) {};
    \node[vertex] (v4) at (-126:1.) {};
    \node[vertex] (v5) at (162:1.) {};
    \draw (v1)--(v2) (v1)--(v3) (v1)--(v4) (v1)--(v5)
    (v2)--(v3) (v2)--(v4) (v2)--(v5)
    (v3)--(v4) (v3)--(v5)
    (v4)--(v5);
  \end{tikzpicture}
}




\subsubsection{Graphs with four or five vertices and maximum degree three}

Tables~\ref{tab:small-simple4}--\ref{tab:small-simple5} give the
complete Hilbert profile of each displayed graph on the cubic parameter
plane (Remark~\ref{rem:small-complete-profiles}). Every
$f=u+\alpha u^2+\beta u^3$ has one of the five listed series: the first
two columns cover $\alpha=0$ (after scaling, $f=u$ or $f=u+u^3$), the
third is the generic value for $\alpha\ne0$, and the last two are the
values on the parabolas $\beta=2\alpha^2$ (inverse quadratics) and
$3\beta=2\alpha^2$ (exponential jets). Entries that differ from the
generic column are in bold; no other value of $\beta/\alpha^2$ is
exceptional for these graphs. Each entry lists $(h_0,\ldots,h_s)$ for
$\hs_G^f=\sum h_it^i$.

\begin{table}[p]
\centering\footnotesize
\setlength{\tabcolsep}{3pt}
\renewcommand{\arraystretch}{1.18}
\begin{tabular}{@{}>{\centering\arraybackslash}m{1.75cm}*{5}{>{\centering\arraybackslash}m{2.85cm}}@{}}
\toprule
$G$ & $u$ & $u+u^3$ & $u+u^2$ & $u+u^2+2u^3$ & $u+\frac{u^2}{2}+\frac{u^3}{6}$ \\ \midrule

\GDfo{.52} & \begin{tabular}{@{}c@{}}$(1,3,3,1)$\end{tabular} & \begin{tabular}{@{}c@{}}$(1,4,3)$\end{tabular} & \begin{tabular}{@{}c@{}}$(1,4,3)$\end{tabular} & \begin{tabular}{@{}c@{}}$(1,4,3)$\end{tabular} & \begin{tabular}{@{}c@{}}$(1,4,3)$\end{tabular} \\[3pt]

\GPaw{.52} & \begin{tabular}{@{}c@{}}$(1,3,5,4,1)$\end{tabular} & \begin{tabular}{@{}c@{}}$(1,4,6,3)$\end{tabular} & \begin{tabular}{@{}c@{}}$(1,4,7,2)$\end{tabular} & \begin{tabular}{@{}c@{}}$(1,4,7,2)$\end{tabular} & \begin{tabular}{@{}c@{}}$(1,4,7,2)$\end{tabular} \\[3pt]

\GDia{.52} & \begin{tabular}{@{}c@{}}$(1,3,6,8,5,1)$\end{tabular} & \begin{tabular}{@{}c@{}}$(1,4,9,8,2)$\end{tabular} & \begin{tabular}{@{}c@{}}$(1,4,10,9)$\end{tabular} & \begin{tabular}{@{}c@{}}$\mathbf{(1,4,9,10)}$\end{tabular} & \begin{tabular}{@{}c@{}}$(1,4,10,9)$\end{tabular} \\[3pt]

\GKfo{.52} & \begin{tabular}{@{}c@{}}$(1,3,6,10,11,$\\$6,1)$\end{tabular} & \begin{tabular}{@{}c@{}}$(1,4,9,12,8,$\\$3,1)$\end{tabular} & \begin{tabular}{@{}c@{}}$(1,4,10,14,5,$\\$3,1)$\end{tabular} & \begin{tabular}{@{}c@{}}$\mathbf{(1,4,9,15,5,}$\\$\mathbf{3,1)}$\end{tabular} & \begin{tabular}{@{}c@{}}$(1,4,10,14,5,$\\$3,1)$\end{tabular} \\[3pt]

\bottomrule\end{tabular}
\caption{Complete cubic Hilbert profiles of the connected simple graphs with four vertices and $\Delta=3$.}\label{tab:small-simple4}
\end{table}

\begin{table}[p]
\centering\footnotesize
\setlength{\tabcolsep}{3pt}
\renewcommand{\arraystretch}{1.18}
\begin{tabular}{@{}>{\centering\arraybackslash}m{1.75cm}*{5}{>{\centering\arraybackslash}m{2.85cm}}@{}}
\toprule
$G$ & $u$ & $u+u^3$ & $u+u^2$ & $u+u^2+2u^3$ & $u+\frac{u^2}{2}+\frac{u^3}{6}$ \\ \midrule

\Gafa{.52} & \begin{tabular}{@{}c@{}}$(1,3,4,3,1)$\end{tabular} & \begin{tabular}{@{}c@{}}$(1,4,5,2)$\end{tabular} & \begin{tabular}{@{}c@{}}$(1,4,7)$\end{tabular} & \begin{tabular}{@{}c@{}}$(1,4,7)$\end{tabular} & \begin{tabular}{@{}c@{}}$(1,4,7)$\end{tabular} \\[3pt]

\Gafb{.52} & \begin{tabular}{@{}c@{}}$(1,3,4,3,1)$\end{tabular} & \begin{tabular}{@{}c@{}}$(1,4,5,2)$\end{tabular} & \begin{tabular}{@{}c@{}}$(1,4,7)$\end{tabular} & \begin{tabular}{@{}c@{}}$(1,4,7)$\end{tabular} & \begin{tabular}{@{}c@{}}$(1,4,7)$\end{tabular} \\[3pt]

\Gafc{.52} & \begin{tabular}{@{}c@{}}$(1,3,5,5,3,1)$\end{tabular} & \begin{tabular}{@{}c@{}}$(1,4,7,5,1)$\end{tabular} & \begin{tabular}{@{}c@{}}$(1,4,9,4)$\end{tabular} & \begin{tabular}{@{}c@{}}$(1,4,9,4)$\end{tabular} & \begin{tabular}{@{}c@{}}$(1,4,9,4)$\end{tabular} \\[3pt]

\Gppw{.52} & \begin{tabular}{@{}c@{}}$(1,3,5,6,4,1)$\end{tabular} & \begin{tabular}{@{}c@{}}$(1,4,8,6,1)$\end{tabular} & \begin{tabular}{@{}c@{}}$(1,4,9,6)$\end{tabular} & \begin{tabular}{@{}c@{}}$\mathbf{(1,4,8,7)}$\end{tabular} & \begin{tabular}{@{}c@{}}$(1,4,9,6)$\end{tabular} \\[3pt]

\Gcfa{.52} & \begin{tabular}{@{}c@{}}$(1,3,6,7,4,1)$\end{tabular} & \begin{tabular}{@{}c@{}}$(1,4,8,7,2)$\end{tabular} & \begin{tabular}{@{}c@{}}$(1,4,10,7)$\end{tabular} & \begin{tabular}{@{}c@{}}$\mathbf{(1,4,9,8)}$\end{tabular} & \begin{tabular}{@{}c@{}}$\mathbf{(1,4,9,8)}$\end{tabular} \\[3pt]

\Gcfb{.52} & \begin{tabular}{@{}c@{}}$(1,3,6,9,8,$\\$4,1)$\end{tabular} & \begin{tabular}{@{}c@{}}$(1,4,10,11,5,1)$\end{tabular} & \begin{tabular}{@{}c@{}}$(1,4,10,15,2)$\end{tabular} & \begin{tabular}{@{}c@{}}$\mathbf{(1,4,9,15,3)}$\end{tabular} & \begin{tabular}{@{}c@{}}$\mathbf{(1,4,10,14,3)}$\end{tabular} \\[3pt]

\bottomrule\end{tabular}
\caption{Complete cubic Hilbert profiles of connected multigraphs with four vertices and $\Delta=3$.}\label{tab:small-multi4}
\end{table}

\begin{table}[p]
\centering\footnotesize
\setlength{\tabcolsep}{3pt}
\renewcommand{\arraystretch}{1.18}
\begin{tabular}{@{}>{\centering\arraybackslash}m{1.75cm}*{5}{>{\centering\arraybackslash}m{2.85cm}}@{}}
\toprule
$G$ & $u$ & $u+u^3$ & $u+u^2$ & $u+u^2+2u^3$ & $u+\frac{u^2}{2}+\frac{u^3}{6}$ \\ \midrule

\GDfi{.52} & \begin{tabular}{@{}c@{}}$(1,4,6,4,1)$\end{tabular} & \begin{tabular}{@{}c@{}}$(1,5,7,3)$\end{tabular} & \begin{tabular}{@{}c@{}}$(1,5,10)$\end{tabular} & \begin{tabular}{@{}c@{}}$(1,5,10)$\end{tabular} & \begin{tabular}{@{}c@{}}$(1,5,10)$\end{tabular} \\[3pt]

\Gban{.52} & \begin{tabular}{@{}c@{}}$(1,4,9,10,5,1)$\end{tabular} & \begin{tabular}{@{}c@{}}$(1,5,11,10,3)$\end{tabular} & \begin{tabular}{@{}c@{}}$(1,5,14,9,1)$\end{tabular} & \begin{tabular}{@{}c@{}}$\mathbf{(1,5,13,10,1)}$\end{tabular} & \begin{tabular}{@{}c@{}}$(1,5,14,9,1)$\end{tabular} \\[3pt]

\Gbul{.52} & \begin{tabular}{@{}c@{}}$(1,4,8,9,5,1)$\end{tabular} & \begin{tabular}{@{}c@{}}$(1,5,12,9,1)$\end{tabular} & \begin{tabular}{@{}c@{}}$(1,5,13,9)$\end{tabular} & \begin{tabular}{@{}c@{}}$\mathbf{(1,5,12,10)}$\end{tabular} & \begin{tabular}{@{}c@{}}$(1,5,13,9)$\end{tabular} \\[3pt]

\Glol{.52} & \begin{tabular}{@{}c@{}}$(1,4,8,9,5,1)$\end{tabular} & \begin{tabular}{@{}c@{}}$(1,5,10,9,3)$\end{tabular} & \begin{tabular}{@{}c@{}}$(1,5,13,8,1)$\end{tabular} & \begin{tabular}{@{}c@{}}$\mathbf{(1,5,12,8,2)}$\end{tabular} & \begin{tabular}{@{}c@{}}$\mathbf{(1,5,13,7,2)}$\end{tabular} \\[3pt]

\Gbip{.52} & \begin{tabular}{@{}c@{}}$(1,4,10,17,15,$\\$6,1)$\end{tabular} & \begin{tabular}{@{}c@{}}$(1,5,15,22,11)$\end{tabular} & \begin{tabular}{@{}c@{}}$(1,5,15,26,7)$\end{tabular} & \begin{tabular}{@{}c@{}}$\mathbf{(1,5,14,25,9)}$\end{tabular} & \begin{tabular}{@{}c@{}}$\mathbf{(1,5,15,24,9)}$\end{tabular} \\[3pt]

\Gkit{.52} & \begin{tabular}{@{}c@{}}$(1,4,9,14,13,$\\$6,1)$\end{tabular} & \begin{tabular}{@{}c@{}}$(1,5,13,19,10)$\end{tabular} & \begin{tabular}{@{}c@{}}$(1,5,14,23,5)$\end{tabular} & \begin{tabular}{@{}c@{}}$\mathbf{(1,5,13,21,8)}$\end{tabular} & \begin{tabular}{@{}c@{}}$\mathbf{(1,5,14,22,6)}$\end{tabular} \\[3pt]

\Ghou{.52} & \begin{tabular}{@{}c@{}}$(1,4,10,16,14,$\\$6,1)$\end{tabular} & \begin{tabular}{@{}c@{}}$(1,5,14,20,11,1)$\end{tabular} & \begin{tabular}{@{}c@{}}$(1,5,15,26,5)$\end{tabular} & \begin{tabular}{@{}c@{}}$\mathbf{(1,5,14,23,9)}$\end{tabular} & \begin{tabular}{@{}c@{}}$\mathbf{(1,5,15,24,7)}$\end{tabular} \\[3pt]

\Gsub{.52} & \begin{tabular}{@{}c@{}}$(1,4,10,19,25,$\\$19,7,1)$\end{tabular} & \begin{tabular}{@{}c@{}}$(1,5,15,29,26,$\\$9,1)$\end{tabular} & \begin{tabular}{@{}c@{}}$(1,5,15,33,28,$\\$3,1)$\end{tabular} & \begin{tabular}{@{}c@{}}$\mathbf{(1,5,14,29,33,}$\\$\mathbf{3,1)}$\end{tabular} & \begin{tabular}{@{}c@{}}$(1,5,15,33,28,$\\$3,1)$\end{tabular} \\[3pt]

\bottomrule\end{tabular}
\caption{Complete cubic Hilbert profiles of the connected simple graphs with five vertices and $\Delta=3$.}\label{tab:small-simple5}
\end{table}


\subsubsection{Simple graphs with five vertices and maximum degree four}

Table~\ref{tab:5v-4d} records the Hilbert functions at nine parameters
for each displayed graph. The columns $f_2,f_5,f_6,f_8,f_9$ lie in the
generic cell of Theorem~\ref{thm:D4-strata} and agree for every graph
in the table, while $f_7=u+u^2+u^3$ lies on $[u^4]f^{-1}=0$ (cell~7) and
differs from them for seven of the eleven graphs. The columns $f_1,f_3,f_4$
have $\alpha=0$. Unlike the previous tables, these columns are not shown
to exhaust the Hilbert strata of these graphs.

\begin{sidewaystable}[p]
\centering\footnotesize
\setlength{\tabcolsep}{3pt}
\renewcommand{\arraystretch}{1.25}
\begin{tabular}{@{}>{\centering\arraybackslash}m{1.65cm}*{9}{>{\centering\arraybackslash}m{2.14cm}}@{}}
\toprule
$G$ & $f_1$ & $f_2$ & $f_3$ & $f_4$ & $f_5$ & $f_6$ & $f_7$ & $f_8$ & $f_9$ \\ \midrule

\Gsta{.45} & \begin{tabular}{@{}c@{}}$(1,5,6,4)$\end{tabular} & \begin{tabular}{@{}c@{}}$(1,5,10)$\end{tabular} & \begin{tabular}{@{}c@{}}$(1,5,7,3)$\end{tabular} & \begin{tabular}{@{}c@{}}$(1,5,7,3)$\end{tabular} & \begin{tabular}{@{}c@{}}$(1,5,10)$\end{tabular} & \begin{tabular}{@{}c@{}}$(1,5,10)$\end{tabular} & \begin{tabular}{@{}c@{}}$(1,5,10)$\end{tabular} & \begin{tabular}{@{}c@{}}$(1,5,10)$\end{tabular} & \begin{tabular}{@{}c@{}}$(1,5,10)$\end{tabular} \\[2pt]

\Gcri{.45} & \begin{tabular}{@{}c@{}}$(1,5,9,9,4)$\end{tabular} & \begin{tabular}{@{}c@{}}$(1,5,13,8,1)$\end{tabular} & \begin{tabular}{@{}c@{}}$(1,5,10,9,3)$\end{tabular} & \begin{tabular}{@{}c@{}}$(1,5,10,9,3)$\end{tabular} & \begin{tabular}{@{}c@{}}$(1,5,13,8,1)$\end{tabular} & \begin{tabular}{@{}c@{}}$(1,5,13,8,1)$\end{tabular} & \begin{tabular}{@{}c@{}}$(1,5,13,8,1)$\end{tabular} & \begin{tabular}{@{}c@{}}$(1,5,13,8,1)$\end{tabular} & \begin{tabular}{@{}c@{}}$(1,5,13,8,1)$\end{tabular} \\[2pt]

\Gbut{.45} & \begin{tabular}{@{}c@{}}$(1,5,12,15,12,$\\$4)$\end{tabular} & \begin{tabular}{@{}c@{}}$(1,5,15,21,7)$\end{tabular} & \begin{tabular}{@{}c@{}}$(1,5,13,16,11,$\\$3)$\end{tabular} & \begin{tabular}{@{}c@{}}$(1,5,13,16,11,$\\$3)$\end{tabular} & \begin{tabular}{@{}c@{}}$(1,5,15,21,7)$\end{tabular} & \begin{tabular}{@{}c@{}}$(1,5,15,21,7)$\end{tabular} & \begin{tabular}{@{}c@{}}$(1,5,15,21,7)$\end{tabular} & \begin{tabular}{@{}c@{}}$(1,5,15,21,7)$\end{tabular} & \begin{tabular}{@{}c@{}}$(1,5,15,21,7)$\end{tabular} \\[2pt]

\Gdar{.45} & \begin{tabular}{@{}c@{}}$(1,5,11,15,12,$\\$4)$\end{tabular} & \begin{tabular}{@{}c@{}}$(1,5,14,25,3)$\end{tabular} & \begin{tabular}{@{}c@{}}$(1,5,14,19,9)$\end{tabular} & \begin{tabular}{@{}c@{}}$(1,5,14,20,8)$\end{tabular} & \begin{tabular}{@{}c@{}}$(1,5,14,25,3)$\end{tabular} & \begin{tabular}{@{}c@{}}$(1,5,14,25,3)$\end{tabular} & \begin{tabular}{@{}c@{}}$(1,5,14,25,3)$\end{tabular} & \begin{tabular}{@{}c@{}}$(1,5,14,25,3)$\end{tabular} & \begin{tabular}{@{}c@{}}$(1,5,14,25,3)$\end{tabular} \\[2pt]

\Ggem{.45} & \begin{tabular}{@{}c@{}}$(1,5,13,21,23,$\\$15,4)$\end{tabular} & \begin{tabular}{@{}c@{}}$(1,5,15,33,28)$\end{tabular} & \begin{tabular}{@{}c@{}}$(1,5,15,28,25,$\\$8)$\end{tabular} & \begin{tabular}{@{}c@{}}$(1,5,15,29,27,$\\$5)$\end{tabular} & \begin{tabular}{@{}c@{}}$(1,5,15,33,28)$\end{tabular} & \begin{tabular}{@{}c@{}}$(1,5,15,33,28)$\end{tabular} & \begin{tabular}{@{}c@{}}$(1,5,15,32,29)$\end{tabular} & \begin{tabular}{@{}c@{}}$(1,5,15,33,28)$\end{tabular} & \begin{tabular}{@{}c@{}}$(1,5,15,33,28)$\end{tabular} \\[2pt]

\GKKK{.45} & \begin{tabular}{@{}c@{}}$(1,5,14,23,23,$\\$14,1)$\end{tabular} & \begin{tabular}{@{}c@{}}$(1,5,15,32,28)$\end{tabular} & \begin{tabular}{@{}c@{}}$(1,5,15,26,25,$\\$9)$\end{tabular} & \begin{tabular}{@{}c@{}}$(1,5,15,27,25,$\\$8)$\end{tabular} & \begin{tabular}{@{}c@{}}$(1,5,15,32,28)$\end{tabular} & \begin{tabular}{@{}c@{}}$(1,5,15,32,28)$\end{tabular} & \begin{tabular}{@{}c@{}}$(1,5,15,31,29)$\end{tabular} & \begin{tabular}{@{}c@{}}$(1,5,15,32,28)$\end{tabular} & \begin{tabular}{@{}c@{}}$(1,5,15,32,28)$\end{tabular} \\[2pt]

\Gllo{.45} & \begin{tabular}{@{}c@{}}$(1,5,11,19,21,$\\$15,4)$\end{tabular} & \begin{tabular}{@{}c@{}}$(1,5,14,29,21,$\\$5,1)$\end{tabular} & \begin{tabular}{@{}c@{}}$(1,5,14,24,22,$\\$9,1)$\end{tabular} & \begin{tabular}{@{}c@{}}$(1,5,14,25,23,$\\$7,1)$\end{tabular} & \begin{tabular}{@{}c@{}}$(1,5,14,29,21,$\\$5,1)$\end{tabular} & \begin{tabular}{@{}c@{}}$(1,5,14,29,21,$\\$5,1)$\end{tabular} & \begin{tabular}{@{}c@{}}$(1,5,14,28,22,$\\$5,1)$\end{tabular} & \begin{tabular}{@{}c@{}}$(1,5,14,29,21,$\\$5,1)$\end{tabular} & \begin{tabular}{@{}c@{}}$(1,5,14,29,21,$\\$5,1)$\end{tabular} \\[2pt]

\Gwhe{.45} & \begin{tabular}{@{}c@{}}$(1,5,13,26,34,$\\$33,18,4)$\end{tabular} & \begin{tabular}{@{}c@{}}$(1,5,15,35,56,$\\$18,4)$\end{tabular} & \begin{tabular}{@{}c@{}}$(1,5,15,33,45,$\\$29,6)$\end{tabular} & \begin{tabular}{@{}c@{}}$(1,5,15,34,47,$\\$28,4)$\end{tabular} & \begin{tabular}{@{}c@{}}$(1,5,15,35,56,$\\$18,4)$\end{tabular} & \begin{tabular}{@{}c@{}}$(1,5,15,35,56,$\\$18,4)$\end{tabular} & \begin{tabular}{@{}c@{}}$(1,5,15,34,56,$\\$19,4)$\end{tabular} & \begin{tabular}{@{}c@{}}$(1,5,15,35,56,$\\$18,4)$\end{tabular} & \begin{tabular}{@{}c@{}}$(1,5,15,35,56,$\\$18,4)$\end{tabular} \\[2pt]

\Ghox{.45} & \begin{tabular}{@{}c@{}}$(1,5,14,27,35,$\\$30,15,1)$\end{tabular} & \begin{tabular}{@{}c@{}}$(1,5,15,34,53,$\\$19,1)$\end{tabular} & \begin{tabular}{@{}c@{}}$(1,5,15,31,43,$\\$28,5)$\end{tabular} & \begin{tabular}{@{}c@{}}$(1,5,15,32,47,$\\$25,3)$\end{tabular} & \begin{tabular}{@{}c@{}}$(1,5,15,34,53,$\\$19,1)$\end{tabular} & \begin{tabular}{@{}c@{}}$(1,5,15,34,53,$\\$19,1)$\end{tabular} & \begin{tabular}{@{}c@{}}$(1,5,15,33,52,$\\$21,1)$\end{tabular} & \begin{tabular}{@{}c@{}}$(1,5,15,34,53,$\\$19,1)$\end{tabular} & \begin{tabular}{@{}c@{}}$(1,5,15,34,53,$\\$19,1)$\end{tabular} \\[2pt]

\GKme{.45} & \begin{tabular}{@{}c@{}}$(1,5,14,30,48,$\\$50,37,13)$\end{tabular} & \begin{tabular}{@{}c@{}}$(1,5,15,35,65,$\\$64,13)$\end{tabular} & \begin{tabular}{@{}c@{}}$(1,5,15,33,55,$\\$59,28,2)$\end{tabular} & \begin{tabular}{@{}c@{}}$(1,5,15,34,58,$\\$64,20,1)$\end{tabular} & \begin{tabular}{@{}c@{}}$(1,5,15,35,65,$\\$64,13)$\end{tabular} & \begin{tabular}{@{}c@{}}$(1,5,15,35,65,$\\$64,13)$\end{tabular} & \begin{tabular}{@{}c@{}}$(1,5,15,34,62,$\\$68,13)$\end{tabular} & \begin{tabular}{@{}c@{}}$(1,5,15,35,65,$\\$64,13)$\end{tabular} & \begin{tabular}{@{}c@{}}$(1,5,15,35,65,$\\$64,13)$\end{tabular} \\[2pt]

\GKfi{.45} & \begin{tabular}{@{}c@{}}$(1,5,14,30,53,$\\$73,60,41,9,4,$\\$1)$\end{tabular} & \begin{tabular}{@{}c@{}}$(1,5,15,35,67,$\\$91,48,15,9,4,$\\$1)$\end{tabular} & \begin{tabular}{@{}c@{}}$(1,5,15,33,60,$\\$76,60,27,9,4,$\\$1)$\end{tabular} & \begin{tabular}{@{}c@{}}$(1,5,15,34,63,$\\$82,56,21,9,4,$\\$1)$\end{tabular} & \begin{tabular}{@{}c@{}}$(1,5,15,35,67,$\\$91,48,15,9,4,$\\$1)$\end{tabular} & \begin{tabular}{@{}c@{}}$(1,5,15,35,67,$\\$91,48,15,9,4,$\\$1)$\end{tabular} & \begin{tabular}{@{}c@{}}$(1,5,15,34,64,$\\$90,53,15,9,4,$\\$1)$\end{tabular} & \begin{tabular}{@{}c@{}}$(1,5,15,35,67,$\\$91,48,15,9,4,$\\$1)$\end{tabular} & \begin{tabular}{@{}c@{}}$(1,5,15,35,67,$\\$91,48,15,9,4,$\\$1)$\end{tabular} \\[2pt]

\bottomrule\end{tabular}
\medskip
\begin{gather*}
f_1=u+u^4,\quad f_2=u+u^2,\quad f_3=u+u^3,\quad
f_4=u+u^3+u^4,\quad f_5=u+u^2+u^4,\\
f_6=u+\tfrac12u^2+\tfrac13u^3,\quad f_7=u+u^2+u^3,\quad
f_8=u+\tfrac12u^2+\tfrac13u^3+\tfrac14u^4,\quad
f_9=u+u^2+u^3+u^4.
\end{gather*}
\caption{Exact Hilbert functions for connected simple graphs with five
vertices and $\Delta=4$, for the nine indicated generating functions.}
\label{tab:5v-4d}
\end{sidewaystable}
\clearpage


\end{document}
